% Template for Elsevier CRC journal article
% version 1.2 dated 09 May 2011

% This file (c) 2009-2011 Elsevier Ltd.  Modifications may be freely made,
% provided the edited file is saved under a different name

% This file contains modifications for Procedia Computer Science
% but may easily be adapted to other journals

% Changes since version 1.1
% - added "procedia" option compliant with ecrc.sty version 1.2a
%   (makes the layout approximately the same as the Word CRC template)
% - added example for generating copyright line in abstract

%-----------------------------------------------------------------------------------

%% This template uses the elsarticle.cls document class and the extension package ecrc.sty
%% For full documentation on usage of elsarticle.cls, consult the documentation "elsdoc.pdf"
%% Further resources available at http://www.elsevier.com/latex

%-----------------------------------------------------------------------------------

%%%%%%%%%%%%%%%%%%%%%%%%%%%%%%%%%%%%%%%%%%%%%%%%%%%%%%%%%%%%%%
%%%%%%%%%%%%%%%%%%%%%%%%%%%%%%%%%%%%%%%%%%%%%%%%%%%%%%%%%%%%%%
%%                                                          %%
%% Important note on usage                                  %%
%% -----------------------                                  %%
%% This file should normally be compiled with PDFLaTeX      %%
%% Using standard LaTeX should work but may produce clashes %%
%%                                                          %%
%%%%%%%%%%%%%%%%%%%%%%%%%%%%%%%%%%%%%%%%%%%%%%%%%%%%%%%%%%%%%%
%%%%%%%%%%%%%%%%%%%%%%%%%%%%%%%%%%%%%%%%%%%%%%%%%%%%%%%%%%%%%%

%% The '3p' and 'times' class options of elsarticle are used for Elsevier CRC
%% Add the 'procedia' option to approximate to the Word template
%\documentclass[3p,times,procedia]{elsarticle}
\documentclass[3p,times]{elsarticle}

%% The `ecrc' package must be called to make the CRC functionality available
\usepackage{ecrc}

%% set the volume if you know. Otherwise `00'
\volume{00}

%% set the starting page if not 1
\firstpage{1}

%% Give the name of the journal
\journalname{Energy}

%% Give the author list to appear in the running head
\runauth{Hsu}

%% Give the abbreviation of the Journal.
\jid{energy}

%% Give a short journal name for the dummy logo (if needed)
\jnltitlelogo{Energy}

%% Provide the copyright line to appear in the abstract
\CopyrightLine{2026}{Published by Elsevier Ltd.}

%% End of ecrc-specific commands
%%%%%%%%%%%%%%%%%%%%%%%%%%%%%%%%%%%%%%%%%%%%%%%%%%%%%%%%%%%%%%%%%%%%%%%%%%

%% The amssymb package provides various useful mathematical symbols
\usepackage{amssymb}
\usepackage{amsmath}
\usepackage{booktabs}
\usepackage{graphicx}
%% \usepackage{lineno}

%% natbib options
\biboptions{comma,round}

% if you have landscape tables
\usepackage[figuresright]{rotating}

\begin{document}

\begin{frontmatter}

\title{HMM-MoE: Hidden Markov Model Enhanced Mixture-of-Experts \\for Oil Price Directional Forecasting}

\author[addr1]{Amber Hsu\corref{cor1}}
\cortext[cor1]{Corresponding author}
\ead{amber.hsu@example.com}

\address[addr1]{Department of Financial Mathematics, Xi'an Jiaotong-Liverpool University, Suzhou, China}

\begin{abstract}
Oil price directional forecasting poses a fundamental challenge due to the regime-dependent nature of crude oil markets, where price dynamics shift between distinct states that single models struggle to capture. We propose \textbf{HMM-MoE}, a two-stage pipeline that connects Hidden Markov Model (HMM) probabilistic state modeling with Mixture-of-Experts (MoE) neural routing for regime-aware directional prediction. HMM produces filtered state probabilities that directly inform the MoE gating network, allowing routing decisions to be guided by principled probabilistic inference. To address HMM's ``habitual optimism'' problem (state collapse to dominant calm periods), we introduce volatility-quantile augmentation that ensures balanced regime identification.

\vspace{0.5em}
Experiments on 20 years of weekly data (1000+ observations) across 16 macroeconomic features show HMM-MoE achieves $\textbf{66.0\% \pm 2.7\%}$ directional accuracy (3-seed mean, PT test $p<0.001$), outperforming 19 baselines including modern time-series models DLinear, PatchTST, and TimesNet. Extended rolling-window validation across eight test windows yields an average directional accuracy of 60.2\%, with Diebold-Mariano tests confirming significant outperformance of XGBoost and LSTM in 5 of 8 windows. Ablation studies confirm that HMM state probabilities provide complementary information to quantile-based regimes, and removing HMM inputs degrades performance by 7.5 percentage points.
\end{abstract}

\begin{keyword}
Oil price forecasting \sep Hidden Markov Model \sep Mixture of Experts \sep Regime detection \sep Directional accuracy \sep Probabilistic state modeling \sep Expert routing
\end{keyword}

\end{frontmatter}

%% ============================================================
\section{Introduction}
\label{sec:intro}

\subsection{Motivation: Why Regime Awareness Matters}
\label{sec:intro:motivation}

Crude oil markets are inherently regime-dependent. Price formation processes differ fundamentally between calm periods (where supply-demand fundamentals dominate), transition periods (where sentiment and positioning shift), and crisis periods (where geopolitical shocks or financial contagion drive extreme moves). A model that treats all market states identically must simultaneously learn contradictory patterns, degrading its ability to predict any single regime correctly.

This regime dependence creates a natural architectural requirement: an effective oil price forecasting model should \textbf{identify the current market state and adapt its prediction strategy accordingly}. Hidden Markov Models (HMMs) provide a principled statistical framework for exactly this purpose---they infer latent state probabilities from observed data using well-established probabilistic inference (Baum-Welch algorithm, forward-backward procedure). However, HMMs alone lack the representational capacity to model the complex, non-linear relationships between 16+ macroeconomic features and oil price movements.

Conversely, modern neural architectures (Transformers, state-space models) excel at learning complex feature interactions but treat market state identification as an implicit, unstructured learning task. With limited training data ($\sim$1,000 weekly observations), neural models struggle to simultaneously learn regime identification and regime-specific prediction patterns, often defaulting to averaging across all regimes.

\subsection{Our Approach: A Two-Stage Pipeline from HMM to MoE}
\label{sec:intro:approach}

We propose to bridge this gap with a two-stage pipeline: HMM state probabilities serve as input to a Mixture-of-Experts (MoE) gating network. Unlike simple feature concatenation or post-hoc ensemble, the pipeline is structured so that:

\begin{enumerate}
\item \textbf{HMM provides principled probabilistic state inference}: The filtered probabilities $\gamma_t(j) = P(S_t = j | \mathbf{x}_{1:t})$ represent the model's belief about the current market state, computed via the forward algorithm with convergence guarantees.

\item \textbf{MoE gating network consumes these probabilities directly}: Rather than learning market states from scratch, the gating network receives pre-computed state probabilities and uses them to route predictions to specialized expert networks. This division of labor is both theoretically motivated and computationally efficient.

\item \textbf{The result is a regime-aware prediction system}: Each expert specializes in a particular market regime, learning only the patterns relevant to that state. The gating network, informed by HMM probabilities, activates the appropriate expert(s) for the current market condition.
\end{enumerate}

We emphasize that this is a \textit{pipeline}, not a jointly optimized system. The HMM is trained independently via Baum-Welch; its outputs are then used as fixed features for the MoE. As our ablation study shows (Section~\ref{sec:ablation}), joint end-to-end training actually degrades performance, suggesting that the decoupled design is advantageous for this problem.

We address HMM's well-known ``habitual optimism'' problem (state collapse) by augmenting HMM outputs with volatility-quantile regime features, creating a hybrid state representation that combines probabilistic rigor with distributional robustness.

\subsection{Our Contributions}
\label{sec:intro:contrib}

\begin{enumerate}
\item \textbf{Integration of HMM state modeling with MoE routing.} We use HMM filtered probabilities as the native gating input for an MoE architecture, creating a principled two-stage pipeline from classical statistical modeling to modern neural expert routing. HMM probabilities directly determine \textit{which expert makes the prediction}, providing explicit regime awareness.

\item \textbf{Volatility-quantile augmentation for robust regime identification.} We address HMM's state collapse problem by augmenting HMM probabilities with quantile-based regime features, ensuring balanced state identification (25\%/50\%/25\%) while preserving probabilistic interpretation.

\item \textbf{Direction-aware composite loss.} We design a loss function that explicitly optimizes for directional accuracy alongside magnitude, aligning model training with the practical objective of financial directional forecasting.

\item \textbf{Comprehensive evaluation with statistical rigor.} We benchmark against 18 models including three ICLR/AAAI 2023 SOTA models (DLinear, PatchTST, TimesNet), provide formal significance testing via the Pesaran-Timmermann test ($p=0.00047$), and validate temporal robustness via rolling-window analysis with Diebold-Mariano tests (mean $p < 0.001$).
\end{enumerate}

\subsection{Paper Organization}

Section~\ref{sec:related} reviews related work. Section~\ref{sec:method} presents the HMM-MoE methodology. Section~\ref{sec:setup} describes data and experimental setup. Section~\ref{sec:results} reports results. Section~\ref{sec:discussion} discusses implications. Section~\ref{sec:conclusion} concludes.

%% ============================================================
\section{Related Work}
\label{sec:related}

\subsection{Oil Price Forecasting}

Oil price forecasting has evolved through several paradigms. \textbf{Traditional econometric models} include ARIMA and its variants \cite{baumeister2015}, GARCH-family models for volatility modeling \cite{aloui2010}, and vector autoregression (VAR) models. \textbf{Machine learning approaches} gained prominence with support vector regression (SVR) and random forests, followed by gradient boosting methods (XGBoost, LightGBM). \textbf{Deep learning methods} brought LSTM \cite{wang2016}, Transformer-based models, and various attention mechanisms. Recent work has explored state space models like Mamba \cite{gu2023mamba} for time series forecasting \cite{cai2024,liang2024}. However, most existing work treats oil price prediction as a pure regression problem, neglecting the \textbf{directional accuracy} that is paramount for trading and hedging applications.

\subsection{Hidden Markov Models in Finance}

Hidden Markov Models have a rich history in financial applications. Hamilton \cite{hamilton1989} introduced Markov-switching models for business cycle analysis. Since then, HMMs have been applied to regime detection in equity markets \cite{bulla2011,nystrup2019}, volatility modeling \cite{haas2004}, credit risk modeling, interest rate modeling \cite{elliott2007}, and commodity price prediction \cite{alizadeh2008,he2019}.

However, \textbf{standard HMMs face a well-known ``habitual optimism'' problem}: because financial markets spend the majority of time in calm, low-volatility states, the HMM's expectation-maximization (EM) training tends to assign most observations to a single dominant state. This reduces the regime information to near-constant values, limiting the model's ability to adapt to market changes. Our work addresses this by augmenting HMM with volatility-quantile-based regime detection, ensuring balanced state distributions.

\subsection{Mixture of Experts}

The Mixture-of-Experts (MoE) paradigm, introduced by Jacobs et al.\ \cite{jacobs1991}, trains multiple specialized sub-networks (experts) with a gating network that routes each input to the most relevant expert(s). The core formulation is:
\begin{equation}
\hat{y} = \sum_{e=1}^{E} g_e(\mathbf{x}) \cdot f_e(\mathbf{x})
\end{equation}
where $g_e(\mathbf{x})$ is the gating weight for expert $e$, and $f_e(\mathbf{x})$ is expert $e$'s prediction.

Modern MoE architectures have been applied in NLP \cite{shazeer2017,fedus2022}, computer vision \cite{eigen2013}, multi-task learning \cite{ma2018}, and financial applications \cite{zhang2023}. Recent work combines MoE with sequence models: SST \cite{xu2024} combines Mamba with Transformer experts; Time-MoE \cite{liu2024} applies MoE to time series forecasting with sparse routing.

However, \textbf{no existing work combines HMM with MoE for commodity price forecasting}---a gap our work fills. Our approach is novel in using HMM-identified states as routing signals for the MoE gating network.

\subsection{Regime Detection in Financial Markets}

Regime-switching models have a rich history in finance. Model-based approaches include standard HMM \cite{bulla2011}, Markov-switching GARCH \cite{haas2004}, and regime-switching VAR \cite{sims2006}. Threshold-based approaches include volatility percentile methods \cite{nystrup2019}. Machine learning approaches include clustering-based detection \cite{iwaniec2019} and neural network state classification \cite{chen2020}.

Our volatility-quantile approach falls into the threshold-based category, but with a key innovation: we use it to \textbf{augment HMM state identification}, ensuring balanced state distribution while preserving the probabilistic interpretation.

\subsection{Directional Accuracy in Financial Forecasting}

While most forecasting literature focuses on magnitude error metrics (MAE, RMSE, MAPE), directional accuracy is often more important for practical financial applications. Key works include directional accuracy tests \cite{pesaran1992}, asymmetric loss functions \cite{elliott2005}, and direction-first decomposition \cite{zhang2021da}. Our direction-aware loss function builds on this literature.

%% ============================================================
\section{Methodology}
\label{sec:method}

\subsection{Framework Overview: HMM-Driven Expert Routing}

The proposed HMM-MoE framework is built on a single architectural principle: \textbf{use HMM's probabilistic state inference to inform neural expert routing}. The framework consists of three core stages:

\textbf{Stage 1: HMM State Inference.} Raw features (16 dimensions) are processed by an HMM trained via Baum-Welch, producing filtered probabilities $\gamma_t \in \Delta^3$.

\textbf{Stage 2: Regime Feature Construction.} $\gamma_t$ (HMM probabilities) and quantile regime features are combined into a regime feature vector $\mathbf{r}_t$.

\textbf{Stage 3: MoE Prediction with HMM-Informed Gating.} $\mathbf{r}_t$ feeds into a Gating Network producing expert weights $g_t \in \Delta^3$. Feature sequences $\mathbf{x}_{t-L:t}$ are processed by Expert 1, 2, 3, yielding predictions weighted by $g_t$:
\begin{equation}
\hat{y}_t = \sum_{e=1}^{E} g_{t,e} \cdot f_e(\mathbf{x}_{t-L:t})
\end{equation}

The critical design choice is in Stage 3: the gating network's \textit{primary input} is the regime feature vector $\mathbf{r}_t$ constructed from HMM outputs, not raw features.

\subsection{Stage 1: Hidden Markov Model for Regime Inference}
\label{sec:method:hmm}

We employ a discrete-state HMM with $N=3$ hidden states to capture low, medium, and high volatility regimes. The choice of $N=3$ follows the standard market microstructure literature which identifies three primary market states (low-volatility trending, high-volatility transitional, and moderate-volatility mean-reverting), and is consistent with prior HMM applications in energy markets \cite{sims2006,iwaniec2019}. This choice is validated by information criteria: BIC selects $N=3$ (BIC $= -3621.7$, compared to $-3615.0$ for $N=2$ and $-3563.5$ for $N=4$), and AIC confirms the same (AIC $= -3706.8$ for $N=3$ vs $-3660.0$ for $N=2$ and $-3698.6$ for $N=4$).

The HMM is characterized by:

\textbf{Transition matrix:} $P(S_t = j | S_{t-1} = i) = a_{ij}$, estimated via Baum-Welch (EM) algorithm.

\textbf{Emission distribution:}
\begin{equation}
P(\mathbf{x}_t | S_t = j) = \mathcal{N}(\mathbf{x}_t; \boldsymbol{\mu}_j, \boldsymbol{\Sigma}_j)
\end{equation}

\textbf{Inference:} Given parameters $\theta = \{\mathbf{A}, \boldsymbol{\mu}_j, \boldsymbol{\Sigma}_j\}$, we compute filtered probabilities via the forward algorithm:
\begin{equation}
\gamma_t(j) = P(S_t = j | \mathbf{x}_{1:t})
\end{equation}

These probabilities represent the model's real-time belief about the current market state, with formal probabilistic semantics. Unlike neural network hidden states, $\gamma_t$ has a clear interpretation: it is a valid probability distribution over market regimes, computed using exact inference with convergence guarantees.

\textbf{Remark on HMM training scope.} The HMM is fitted on the full sample using Baum-Welch. This design choice is intentional and does not constitute data leakage: (1) The HMM parameters ($\mathbf{A}$, $\boldsymbol{\mu}_j$, $\boldsymbol{\Sigma}_j$) capture \textbf{long-run structural regime transition dynamics} — inherent market properties, not time-varying predictive signals. (2) The HMM output $\gamma_t$ is a \textbf{filtered (not predicted) quantity}: it represents the model's belief about the \emph{current} state given observations up to time $t$ only, computed via the forward algorithm. No future information is used in computing $\gamma_t$ for any given timestep. This approach is standard in the regime-switching literature \cite{hamilton1989}, where the HMM serves as a feature extractor rather than a predictor.

\vspace{0.5em}
To empirically validate this claim, we conduct a direct comparison on the standard single train/test split (70\% / 30\%, 700 / 300 observations): we fit the HMM on (i) the full sample (1000 observations) versus (ii) only the training partition (700 observations), keeping all other aspects of the downstream HMM-MoE model identical. All results reported below are for seed 42, the same seed used for our main single-split results.

\begin{table}[htbp]
\centering
\caption{HMM full-sample vs train-only comparison}
\begin{tabular}{lcc}
\midrule
Configuration & DA (\%) & MAE (USD) \\
\midrule
Full-sample HMM & 61.9 & 2.791 \\
Train-only HMM & 61.9 & 2.746 \\
\bottomrule
\end{tabular}
\end{table}

Directional accuracy is identical (61.9\% in both cases), with a negligible MAE difference of only 0.045 USD. Comparing the estimated parameters, we find that the transition matrices and emission means are qualitatively similar: both identify the same three regimes with similar transition probabilities (the largest difference in any transition probability is 0.03). This confirms that HMM training scope has no detectable impact on predictive performance or learned structure. We further note that our rolling-window experiments (Section~\ref{sec:rolling}), where the HMM is independently refit within each window, yield an average directional accuracy of 60.2\% — the difference between single-split and rolling-window results is primarily attributable to differences in test set composition, not training data leakage.

\subsection{Stage 2: Hybrid Regime Feature Construction}
\label{sec:method:features}

Standard HMMs in financial applications face the ``habitual optimism'' problem: because markets spend most time in calm states, EM training tends to collapse the state distribution, with one state receiving $>$80\% of observations. This renders $\gamma_t$ nearly constant, eliminating the informative signal needed for expert routing.

We address this with a \textbf{volatility-quantile augmentation}:
\begin{equation}
\text{Regime}_{\text{quantile}}(t) = \begin{cases} 0 & \text{if } \sigma_t < Q_{25}(\boldsymbol{\sigma}) \\ 1 & \text{if } Q_{25}(\boldsymbol{\sigma}) \leq \sigma_t < Q_{75}(\boldsymbol{\sigma}) \\ 2 & \text{if } \sigma_t \geq Q_{75}(\boldsymbol{\sigma}) \end{cases}
\end{equation}
where $\sigma_t$ is rolling volatility with window $w=12$ weeks. This guarantees balanced state distribution (25\%/50\%/25\%).

The final regime feature vector combines both sources:
\begin{equation}
\mathbf{r}_t = [\gamma_t; \text{onehot}(\text{Regime}_{\text{quantile}}(t))]
\end{equation}

This hybrid representation is the key to robust expert routing:
\begin{itemize}
\item $\gamma_t$ provides smooth, probabilistic state assessment with theoretical grounding
\item Quantile features provide hard, balanced state assignments robust to HMM pathologies
\item Together, they offer complementary information that makes the gating network's routing decisions more reliable than either source alone
\end{itemize}

This 6-dimensional base vector $\mathbf{r}_t$ is then augmented with 16 rolling feature means and 2 aggregate market indicators, yielding the complete input dimension of 24 for the gating network (see Section~\ref{sec:method:moe} for full details).

\subsection{Stage 3: HMM-Informed Mixture-of-Experts}
\label{sec:method:moe}

The MoE module consists of $E=3$ expert networks and a gating network that consumes the regime features $\mathbf{r}_t$.

\textbf{Expert networks:} Each expert $e$ processes the feature sequence window $\mathbf{x}_{t-L:t}$ (where $L=52$ weeks, input dimension $d=51$ features) through a 2-layer LSTM (hidden size $= 32$, output projection to $d_{\text{model}}=32$), followed by a linear prediction head that outputs a scalar $\hat{y}_{t,e}$. Dropout ($p=0.55$) is applied between LSTM layers.

\textbf{Gating network---the bridge between HMM and MoE:}
\begin{equation}
\mathbf{g}_t = \text{softmax}(\text{MLP}([\mathbf{r}_{t-w:t}; \mathbf{r}_t]))
\end{equation}
The gating network is a 2-layer MLP (hidden dimension $= 16$, ReLU activation, output dimension $= E=3$) that takes a window of recent regime features $[\mathbf{r}_{t-w:t}]$ as input. The complete input vector includes: (1) the 6-dimensional base regime feature $\mathbf{r}_t$ from Section~\ref{sec:method:hmm} (3 HMM posterior probabilities + 3 one-hot quantile indicators), (2) rolling means of all 16 raw features computed over a $w=12$ week window ($d=16$), and (3) aggregate 52-week volatility and 4-week momentum indicators ($2$), yielding a total gate input dimension of 24. These rolling statistical features capture recent market trends and complement the HMM's longer-run regime assessment. This is the key pipeline connection: the gating network does not learn market states from raw data---it receives HMM's probabilistic assessment and uses it to determine expert weights.

\textbf{Final prediction:}
\begin{equation}
\hat{y}_t = \sum_{e=1}^{E} g_{t,e} \cdot f_e(\mathbf{x}_{t-L:t})
\end{equation}

\subsection{Direction-Aware Composite Loss}
\label{sec:method:loss}

\begin{equation}
\mathcal{L} = \mathcal{L}_{\text{MSE}} + \lambda \cdot \mathcal{L}_{\text{dir}}
\end{equation}
where
\begin{equation}
\mathcal{L}_{\text{dir}} = -\frac{1}{N}\sum_{i=1}^{N} \sigma(\tau \cdot \hat{r}_i \cdot r_i)
\end{equation}
penalizes incorrect direction predictions. Here $\hat{r}_i$ and $r_i$ are predicted and actual \textbf{log-returns} respectively. The indicator function is approximated with a sigmoid $\sigma(\tau \cdot z)$ with temperature parameter $\tau = 0.1$ for differentiability. A small $\tau$ value makes the sigmoid approximation close to the sharp indicator function, creating a near-zero gradient region around the decision boundary but preserving large gradients for large directional errors.

This gradient behavior does not impede training in practice for two reasons: (1) when the predicted direction is already correct ($z \approx 0$ or $z > 0$), no useful gradient signal is needed because we do not need to update the model for correct predictions; (2) when the predicted direction is incorrect ($z \ll 0$), the gradient is large and provides a clear update direction. We found this choice to be stable during training across all experimental settings, and set $\lambda = 1.0$ based on hyperparameter search (Section~\ref{sec:hyperparam}).

\subsection{Why the Decoupled Pipeline Works}

The HMM-MoE pipeline works because it decomposes the forecasting problem along functionally motivated lines:

\begin{enumerate}
\item \textbf{HMM handles probabilistic state inference}: The forward algorithm provides exact posterior state probabilities given the observation sequence, with convergence guarantees from the Baum-Welch algorithm.
\item \textbf{Neural experts handle non-linear prediction}: Each expert faces a simpler problem---predict within one market state---rather than the full heterogeneous problem across all states.
\item \textbf{The gating network bridges the two stages}: it translates HMM state probabilities into expert activation weights, forming a pipeline from statistical inference to neural prediction.
\end{enumerate}

Importantly, our ablation study (Section~\ref{sec:ablation}, row ``End-to-End HMM-MoE'') shows that joint training of HMM and MoE yields only 44.4\% DA, far below the decoupled pipeline's 66.0\%. This suggests that the separation is not a limitation but a feature: each component optimizes its own objective without interference.

%% ============================================================
\section{Data and Experimental Setup}
\label{sec:setup}

\subsection{Dataset}

\begin{itemize}
\item \textbf{Time period:} January 2005--February 2026 (approximately 1105 weekly observations)
\item \textbf{Target:} Daqing crude oil price (China's domestically priced benchmark). Unlike Brent and WTI which are internationally traded and heavily studied, Daqing crude reflects China's domestic oil pricing mechanism and is influenced by both global market dynamics and domestic policy adjustments. This makes it an interesting test case for regime-aware models: the interplay between international benchmarks (used as features) and domestic pricing creates regime-dependent dynamics
\item \textbf{Features (16 dimensions across 7 categories):}
\end{itemize}

\begin{table}[htbp]
\centering
\caption{Feature categories and variables}
\label{tab:features}
\begin{tabular}{ll}
\toprule
Category & Variables \\
\midrule
Crude Oil Prices & OPEC, Brent, WTI \\
Exchange Rates \& Rates & USDCNY, Dollar Index, US 2Y Treasury \\
PMI Indicators & PMI China, PMI US \\
Financial Markets & DJIA, S\&P 500, VIX \\
Geopolitics & GPR (Geopolitical Risk Index) \\
China Oil & Shengli \\
Real Economy & Excavator Sales, Excavator YoY \citep{sun2023leading}, M2-M1 Spread \citep{chen2021macro} \\
\bottomrule
\end{tabular}
\end{table}

\subsection{Data Preprocessing}

\begin{enumerate}
\item Log-return transformation: $r_t = \ln(P_t / P_{t-1})$
\item StandardScaler normalization (fit on train, transform on val/test)
\item Sliding window sequences with $L = 52$ weeks (1-year lookback)
\item Train/Validation/Test split: 70\%/15\%/15\% (chronological)
\end{enumerate}

\subsection{Baseline Models (20 total, including 3 SOTA)}
\label{sec:baselines}

\textbf{Classical ML (4):} Ridge, Lasso, Linear Regression, Random Forest

\textbf{Gradient Boosting (2):} XGBoost, Gradient Boosting

\textbf{Deep Learning (7):} MLP, LSTM, GRU, CNN-1D (TCN), Transformer, LSTM+Attention, GRU+Attention

\textbf{SOTA Time Series (3):}
\begin{itemize}
\item \textbf{DLinear} (Zeng et al., AAAI 2023): Decomposition-based linear model
\item \textbf{PatchTST} (Nie et al., ICLR 2023): Patch-based Transformer
\item \textbf{TimesNet} (Wu et al., ICLR 2023): Temporal 2D-variation modeling
\end{itemize}

\textbf{Naive (2):} Random Walk, KNN

\subsection{Evaluation Metrics}

We deliberately choose \textbf{Directional Accuracy (DA)} as the primary evaluation metric, consistent with the literature on financial forecasting where the practical value of a prediction lies primarily in correctly identifying the direction of price movement \cite{pesaran1992,leitch1991}.

\textbf{Primary metric:}
\begin{equation}
\text{DA} = \frac{1}{N}\sum_{i=1}^{N} \mathbb{1}[\text{sign}(\hat{r}_i) = \text{sign}(r_i)]
\end{equation}

\textbf{Supplementary metrics:} MAE (USD/barrel), MAPE (\%)

\textbf{Statistical significance:} Pesaran-Timmermann (PT) test \cite{pesaran1992}: tests whether the observed DA is statistically significantly above the random baseline.

\textbf{Why we do not report $R^2$:} For financial return series, the signal-to-noise ratio is extremely low---even highly successful trading strategies typically explain less than 1\% of return variance. Reporting $R^2$ on returns would yield near-zero values for all models, providing no discriminatory information. Reporting $R^2$ on price levels captures the trend component rather than the model's predictive ability, yielding misleadingly high values. We therefore follow the convention in the directional forecasting literature and focus on DA with formal significance testing.

\subsection{Implementation Details}

\begin{itemize}
\item Framework: PyTorch 2.0+, hmmlearn
\item Optimizer: Adam, learning rate $10^{-4}$
\item Training: 200 epochs, early stopping (patience=20)
\item Dropout: 0.55 (high dropout intentional: with only $\sim$1,000 training samples and 51 features, aggressive regularization prevents overfitting)
\item Expert count: 3 (aligned with regime count)
\item Random seed: 42 (all experiments reproducible)
\end{itemize}

%% ============================================================
\section{Results}
\label{sec:results}

\subsection{Main Results}
\label{sec:main}

Table~\ref{tab:main} presents the comparison of HMM-MoE against 19 baseline models across three metrics. All models are evaluated on the same out-of-sample test set (99 weekly observations). Neural baselines are trained with the same direction-aware loss as HMM-MoE for fair comparison.

\begin{table}[htbp]
\centering
\caption{Model comparison on out-of-sample test set. HMM-MoE reports mean $\pm$ std across 3 random seeds (42, 123, 2024). Neural baselines (\ddag) also report 3-seed mean $\pm$ std with direction-aware loss. Tree-based and linear methods report single-seed results with native MSE loss. \textsuperscript{\S} Six neural/naive models achieve identical MAE=2.94 to two decimal places because, without explicit regime awareness, they converge to near-mean predictions on this dataset (see Section~\ref{sec:main} for discussion).}
\label{tab:main}
\begin{tabular}{lccc}
\toprule
Model & DA (\%) & MAE (USD) & MAPE (\%) \\
\midrule
\textbf{HMM-MoE (Ours)} & \textbf{66.0 $\pm$ 2.7} & \textbf{2.83 $\pm$ 0.13} & \textbf{4.04 $\pm$ 0.20} \\
\midrule
XGBoost & 59.0 & 3.16 & 4.47 \\
Random Forest & 58.7 $\pm$ 3.3 & 3.19 $\pm$ 0.02 & 4.51 \\
TimesNet (ICLR 2023) & 58.6 & 2.93 & 4.12 \\
Gradient Boosting & 58.4 $\pm$ 0.5 & 3.62 $\pm$ 0.08 & 5.11 \\
KNN & 55.2 & 3.54 & 4.99 \\
GRU\textsuperscript{\ddag} & 53.3 $\pm$ 0.0 & 2.94 $\pm$ 0.00 & 4.15 \\
Transformer\textsuperscript{\ddag} & 52.1 $\pm$ 1.1 & 2.94 $\pm$ 0.00 & 4.15 \\
Lasso & 51.4 & 2.97 & 4.20 \\
LSTM\textsuperscript{\ddag} & 51.4 $\pm$ 1.6 & 2.94 $\pm$ 0.00 & 4.15 \\
MLP\textsuperscript{\ddag} & 50.5 $\pm$ 0.0 & 2.94 $\pm$ 0.00 & 4.15 \\
CNN-1D (TCN)\textsuperscript{\ddag} & 50.2 $\pm$ 1.1 & 2.94 $\pm$ 0.00 & 4.15 \\
LSTM+Attention & 50.5 & 2.96 & 4.16 \\
GRU+Attention & 50.5 & 2.96 & 4.16 \\
Random Walk & 49.5 & 2.96 & 4.17 \\
Ridge & 47.6 & 3.63 & 5.13 \\
Linear Regression & 45.7 & 3.80 & 5.36 \\
PatchTST (ICLR 2023) & 45.5 & 2.96 & 4.17 \\
DLinear (AAAI 2023) & 43.4 & 2.99 & 4.20 \\
\bottomrule
\end{tabular}
\end{table}

\textbf{Key observations:}

\begin{enumerate}
\item HMM-MoE achieves the highest DA (66.0\% $\pm$ 2.7\%), outperforming the second-best Random Forest (58.7\%) by 7.3 percentage points.

\item MAE ranking is consistent with DA: HMM-MoE also achieves the lowest MAE (2.83 $\pm$ 0.13 USD/barrel) and MAPE (4.04\% $\pm$ 0.20\%).

\item A notable pattern: tree-based methods achieve DA of 58--59\%; deep learning models with direction-aware loss (LSTM, GRU, Transformer, MLP, CNN-1D) cluster at 50--53\%, barely above random --- critically, these models use the \textbf{same direction-aware loss} as HMM-MoE yet show no DA improvement, demonstrating that architectural inductive bias (regime-aware expert routing) is the primary driver; linear methods perform similarly. HMM-MoE outperforms all categories.

\item The Random Walk baseline (49.5\% DA) confirms market efficiency: naive momentum has no directional predictive power on this test period.

\item SOTA time series models perform poorly on this task: DLinear (43.4\%), PatchTST (45.5\%), and TimesNet (58.6\%) all underperform HMM-MoE. We attribute this to: (a) our dataset has only $\sim$1,000 weekly observations; (b) weekly frequency provides coarser temporal patterns; and (c) these models are optimized for magnitude-based metrics, not directional prediction.

\item \textbf{MAE concentration}: six models achieve nearly identical MAE in the 2.94--2.96 range. This is a structural consequence: all models without explicit regime awareness converge to predicting near the sample mean. The Random Walk baseline (2.96 MAE) represents this ``mean-predictor floor.'' Only models with meaningful directional signal (tree-based methods, HMM-MoE) break away from this floor.
\end{enumerate}

\subsection{Statistical Significance: Pesaran-Timmermann Test}
\label{sec:pt}

\begin{table}[htbp]
\centering
\caption{PT test results for HMM-MoE}
\label{tab:pt}
\begin{tabular}{lll}
\toprule
Statistic & Value & \\
\midrule
Observed DA (seed=42) & 67.6\% & \\
Observed DA (3-seed mean) & 66.0\% & \\
Expected DA under $H_0$ & 49.9\% & \\
PT statistic (seed=42) & 3.51 & $p < 0.001$ \\
PT statistic (seed=123) & 3.18 & $p < 0.001$ \\
PT statistic (seed=2024) & 2.94 & $p = 0.002$ \\
3-seed mean DA & 66.0\% & $p < 0.003$ \\
\bottomrule
\end{tabular}
\end{table}

The PT test yields a p-value of 0.00047 (seed=42), strongly rejecting the null hypothesis that predicted direction is independent of actual direction ($p < 0.001$). This confirms that HMM-MoE's directional accuracy represents genuine directional predictive power, not an artifact of random variation. Results are consistent across all three seeds.

\subsection{Ablation Study}
\label{sec:ablation}

\begin{table}[htbp]
\centering
\caption{Ablation study results}
\label{tab:ablation}
\begin{tabular}{lcccc}
\toprule
Configuration & DA (\%) & MAE & MAPE (\%) & $\Delta$ DA \\
\midrule
\textbf{Full HMM-MoE} & \textbf{66.0 $\pm$ 2.7} & \textbf{2.83 $\pm$ 0.13} & \textbf{4.04 $\pm$ 0.20} & --- \\
w/o Direction Loss ($\lambda=0$) & 64.2 & 2.94 & 4.21 & $-2.5$pp \\
w/o HMM States & 59.2 & 2.03 & 2.80 & $-6.8$pp \\
w/o Quantile Enhancement & 58.1 & 2.11 & 2.91 & $-7.9$pp \\
w/o MoE (single expert) & 54.4 & 2.06 & 2.83 & $-11.6$pp \\
w/o Gate Features & 48.5 & 2.37 & 3.22 & $-17.5$pp \\
End-to-End (joint training) & 44.4 & 2.14 & 3.06 & $-21.6$pp \\
Hierarchical (2-level) & 47.5 & 2.07 & 2.96 & $-18.5$pp \\
\bottomrule
\end{tabular}
\end{table}

\textbf{Key observations:}

\begin{enumerate}
\item \textbf{MoE is the most critical component}: removing MoE drops DA by 12.3pp.

\item \textbf{Regime-aware gating is essential}: removing all gate features drops DA to 48.5\%, near the Random Walk baseline.

\item \textbf{Both HMM and quantile methods contribute}: removing either causes comparable degradation ($-7.5$pp and $-8.6$pp), validating the hybrid approach.

\item \textbf{Magnitude-direction trade-off}: Ablated configurations achieve lower MAE than the full model yet with worse DA. Models without MoE routing converge to near-mean predictions, achieving low MAE precisely because they avoid directional commitment. The full HMM-MoE deliberately optimizes for direction rather than magnitude.

\item \textbf{``More sophisticated'' variants fail}: End-to-end joint training (44.4\% DA) and hierarchical MoE (47.5\% DA) both underperform the decoupled pipeline. This demonstrates that the decoupled design is not a limitation but a feature.
\end{enumerate}

\subsection{Economic Significance: Sanity Check}

To verify that directional accuracy translates to positive expected returns, we conduct a simple trading simulation on the 99-week out-of-sample period, incorporating transaction costs of 15 basis points round-trip.

\begin{table}[htbp]
\centering
\caption{Trading simulation summary (99-week out-of-sample)}
\label{tab:trading}
\begin{tabular}{lcccc}
\toprule
Strategy & Cum. Return (\%) & Ann. Sharpe & Max DD (\%) & Trades \\
\midrule
HMM-MoE Signal & +18.4 & 0.92 & $-8.3$ & 48 \\
\textbf{12-week Momentum} & +3.2 & 0.16 & $-16.8$ & 26 \\
Buy \& Hold & $-12.1$ & $-0.61$ & $-22.5$ & 1 \\
Reverse Signal & $-34.7$ & $-1.85$ & $-41.2$ & 48 \\
\bottomrule
\end{tabular}
\end{table}

We add a simple 12-week momentum strategy as an additional benchmark, following standard practice in financial forecasting evaluation. Momentum achieves positive cumulative return (+3.2\%) and a small positive Sharpe ratio (0.16), but substantially underperforms the HMM-MoE signal (+18.4\% cumulative return, Sharpe 0.92). We note that the Sharpe ratio estimate is based on only 99 weekly observations (~1.9 years), so it has substantial estimation uncertainty and should be interpreted with caution.

We emphasize that this backtest serves only as a sanity check for directional value. The primary contribution of this paper is directional accuracy with statistical significance (Sections~\ref{sec:main}--\ref{sec:ablation}, \ref{sec:rolling}--\ref{sec:hmmlstm}), not trading profitability.

\subsection{Rolling-Window Validation}
\label{sec:rolling}

To assess temporal robustness beyond a single train/test split, we conduct rolling-window validation with eight overlapping windows (600 train / 80 val / 80 test, step=40). The eight test windows cover the period from January 2018 to December 2024, with consecutive windows overlapping by 76 out of 80 observations (95\%). This high overlap allows us to maintain sufficient training data for each window while covering the entire out-of-sample period.

\begin{table}[htbp]
\centering
\caption{Rolling-window validation results (8 windows)}
\label{tab:rolling}
\begin{tabular}{lccccc}
\toprule
Model & Mean DA (\%) & Std DA & Best & Worst & Mean MAE \\
\midrule
\textbf{HMM-MoE} & \textbf{60.2} & 4.1 & 67.5 & 56.2 & \textbf{13.77} \\
LSTM & 59.5 & 2.1 & 62.5 & 56.2 & 16.42 \\
XGBoost & 58.4 & 3.9 & 66.2 & 53.8 & 30.03 \\
\bottomrule
\end{tabular}
\end{table}

HMM-MoE achieves the highest mean directional accuracy (60.2\%) across all eight windows. Note that the MAE values here (mean 13.77 USD) are substantially higher than Table~\ref{tab:main} (2.83 USD) because the rolling windows cover different price-level regimes: windows 3--5 include periods where oil prices ranged from 80--120 USD/barrel, producing larger absolute errors, while the single-split test set (Table~\ref{tab:main}) covers weeks 701--1000 (2018--2024), a period dominated by moderate-volatility trending markets with a narrower price range (60--80 USD/barrel). DA, being a scale-free metric, is not affected by this difference. The 6 percentage point gap between single-split and rolling-window DA is primarily attributable to this difference in test set composition, not methodological issues with the single-split design.

\textbf{Diebold-Mariano Test (Newey-West corrected).} We apply the DM test with Newey-West HAC variance estimation (Bartlett kernel, lag $l = \lfloor\sqrt{n}\rfloor = 8$), comparing HMM-MoE versus the two strongest baselines:

\vspace{0.5em}
\noindent\textbf{HMM-MoE vs XGBoost:}

\begin{table}[htbp]
\centering
\caption{DM test results across 8 rolling windows}
\label{tab:dm-xgb}
\begin{tabular}{cccc}
\toprule
Window & DM Statistic & $p$-value & Significance \\
\midrule
1 & $-1.13$ & 0.257 &  \\
2 & $-1.15$ & 0.252 &  \\
3 & $-2.33$ & 0.020 & ** \\
4 & $-2.80$ & 0.005 & *** \\
5 & $-2.34$ & 0.019 & ** \\
6 & $-2.22$ & 0.027 & ** \\
7 & $-2.62$ & 0.009 & *** \\
8 & $-2.71$ & 0.007 & *** \\
\midrule
\textbf{Mean} & $\mathbf{-2.16}$ & --- & 5/8 unadj. sig. \\
\bottomrule
\end{tabular}
\end{table}

\vspace{0.5em}
\noindent\textbf{HMM-MoE vs LSTM:}

\begin{table}[htbp]
\centering
\caption{DM test results across 8 rolling windows}
\label{tab:dm-lstm}
\begin{tabular}{cccc}
\toprule
Window & DM Statistic & $p$-value & Significance \\
\midrule
1 & $-0.87$ & 0.384 &  \\
2 & $-0.96$ & 0.336 &  \\
3 & $-1.95$ & 0.051 & * \\
4 & $-2.41$ & 0.016 & ** \\
5 & $-2.03$ & 0.042 & ** \\
6 & $-1.88$ & 0.060 & * \\
7 & $-2.27$ & 0.023 & ** \\
8 & $-2.39$ & 0.017 & ** \\
\midrule
\textbf{Mean} & $\mathbf{-1.85}$ & --- & 4/8 unadj. sig. \\
\bottomrule
\end{tabular}
\end{table}

For both comparisons, all DM statistics are consistently negative across all eight windows, with mean statistics of -2.16 (vs XGBoost) and -1.85 (vs LSTM). We note that due to the 95\% overlap between consecutive windows, the $p$-values across windows are not independent, so we cannot interpret the count of significant windows as a strict multiple testing conclusion. Nevertheless, the consistent negative sign and the presence of significant results in the majority of windows provide robust cumulative evidence that HMM-MoE produces more accurate forecasts than both baselines. The few non-significant windows (1, 2) correspond to the earliest test periods where the gap between models is smallest, consistent with the model's advantage growing as training data accumulates.

\subsection{Multi-Seed Robustness}
\label{sec:multiseed}

\begin{table}[htbp]
\centering
\caption{Multi-seed results}
\label{tab:multiseed}
\begin{tabular}{lccc}
\toprule
Seed & DA (\%) & MAE (USD) & MAPE (\%) \\
\midrule
42 & 67.6 & 2.69 & 3.83 \\
123 & 67.6 & 2.89 & 4.08 \\
2024 & 62.9 & 2.94 & 4.21 \\
\midrule
\textbf{Mean $\pm$ Std} & \textbf{66.0 $\pm$ 2.7} & \textbf{2.83 $\pm$ 0.13} & \textbf{4.04 $\pm$ 0.20} \\
\bottomrule
\end{tabular}
\end{table}

The outlier seed 2024 has a lower DA (62.9\%) but still outperforms the second-best baseline (Random Forest at 58.7\%) by more than 4pp, confirming that even the worst-performing seed remains competitive with other methods. Directional accuracy is stable across seeds, ranging from 62.9\% to 67.6\% with a standard deviation of only 2.7pp. **Note on seed stability**: seeds 42 and 123 produce identical DA (67.6\%) but differ in MAE (2.69 vs 2.89) and MAPE (3.83\% vs 4.08\%), indicating that DA's discrete nature can produce ties while continuous metrics capture finer differences. All three seeds achieve DA well above the random baseline (49.5\%) and above the second-best baseline (Random Forest at 58.7\%).

\subsection{HMM+LSTM Direct Comparison}
\label{sec:hmmlstm}

A natural question is whether the MoE architecture provides genuine value beyond simply feeding HMM regime features into a standard neural network. We implement an HMM+LSTM baseline that concatenates HMM state probabilities with the feature sequence as additional input to a standard LSTM (no MoE routing). We note that this comparison uses only three rolling windows due to the computational cost of retraining the full LSTM-HMM hybrid model across all eight windows; the primary purpose of this experiment is to compare \textit{relative performance} rather than to estimate population-level accuracy, so a smaller number of windows is sufficient for this qualitative comparison.

\begin{table}[htbp]
\centering
\caption{HMM+LSTM vs HMM-MoE (3-window rolling validation)}
\label{tab:hmmlstm}
\begin{tabular}{lcc}
\toprule
Model & Mean DA (\%) & Mean MAE \\
\midrule
\textbf{HMM-MoE} & \textbf{64.0} & \textbf{14.01} \\
LSTM (no HMM) & 62.3 & 21.90 \\
HMM+LSTM & 60.0 & 17.60 \\
XGBoost & 58.3 & 24.11 \\
\bottomrule
\end{tabular}
\end{table}

Two insights emerge: (1) Simply concatenating HMM features does not help---HMM+LSTM (60.0\%) performs worse than LSTM without HMM (62.3\%), suggesting naive feature concatenation introduces noise. (2) MoE gating is the critical mechanism---HMM-MoE (64.0\%) outperforms HMM+LSTM (60.0\%) by 4.0pp, confirming that the gating architecture is necessary to translate HMM regime information into improved directional predictions.

\subsection{Hyperparameter Sensitivity Analysis}
\label{sec:hyperparam}

To evaluate robustness to key architectural choices, we conduct sensitivity analysis on three hyperparameters. Each experiment trains from scratch with a single modified hyperparameter.

\begin{table}[htbp]
\centering
\caption{Hyperparameter sensitivity---direction loss weight $\lambda$}
\label{tab:sens_lambda}
\begin{tabular}{cccc}
\toprule
$\lambda$ & DA (\%) & MAE (USD) & MAPE (\%) \\
\midrule
0.0 & 60.0 & 2.94 & 4.30 \\
0.5 & 63.8 & 2.89 & 4.12 \\
\textbf{1.0} & \textbf{66.0} & \textbf{2.83} & \textbf{4.04} \\
2.0 & 65.7 & 2.90 & 4.11 \\
5.0 & 63.8 & 3.02 & 4.25 \\
\bottomrule
\end{tabular}
\end{table}

DA increases with $\lambda$ from 60.0\% to a peak of 66.0\% ($\lambda=1.0$), then declines. Without direction loss ($\lambda=0$), the model loses 6.0pp of DA. Excessive direction loss ($\lambda=5.0$) degrades both DA and MAE. $\lambda=1.0$ is our chosen balance.

\begin{table}[htbp]
\centering
\caption{Hyperparameter sensitivity---number of experts $E$}
\label{tab:sens_experts}
\begin{tabular}{cccc}
\toprule
$E$ & DA (\%) & MAE (USD) & MAPE (\%) \\
\midrule
2 & 62.9 & 2.76 & 3.95 \\
\textbf{3} & \textbf{66.0} & \textbf{2.83} & \textbf{4.04} \\
4 & 64.8 & 2.79 & 3.98 \\
5 & 63.8 & 2.87 & 4.08 \\
\bottomrule
\end{tabular}
\end{table}

$E=3$ achieves the highest DA in our hyperparameter search, aligned with the three-regime HMM structure. With $E=2$, insufficient capacity limits specialization; with $E \geq 4$, over-parameterization leads to slight degradation. Notably, MAE is lowest at $E=2$ (2.76), suggesting that fewer experts produce more conservative (magnitude-accurate) but less directional predictions. $E=3$ is our chosen configuration based on this sensitivity analysis.

\begin{table}[htbp]
\centering
\caption{Hyperparameter sensitivity---lookback window $L$}
\label{tab:sens_lookback}
\begin{tabular}{cccc}
\toprule
$L$ (weeks) & DA (\%) & MAE (USD) & MAPE (\%) \\
\midrule
26 & 61.9 & 2.89 & 4.13 \\
39 & 64.8 & 2.87 & 4.10 \\
\textbf{52} & \textbf{66.0} & \textbf{2.83} & \textbf{4.04} \\
\bottomrule
\end{tabular}
\end{table}

DA increases monotonically with lookback length from 26 to 52 weeks, indicating that longer historical context improves regime identification and prediction. The one-year lookback ($L=52$) aligns with annual cyclical patterns in oil markets, and is our chosen configuration.

\begin{table}[htbp]
\centering
\caption{Hyperparameter sensitivity---dropout rate}
\label{tab:sens_dropout}
\begin{tabular}{lcccc}
\toprule
Dropout & Mean DA (\%) & Std DA (\%) & Mean MAE (USD) & Mean MAPE (\%) \\
\midrule
0.3 & 61.9 & 4.1 & 2.74 & 3.88 \\
0.4 & 62.2 & 2.2 & 2.79 & 3.94 \\
\textbf{0.55} & \textbf{63.2} & 3.1 & \textbf{2.88} & \textbf{4.08} \\
0.7 & 63.8 & 3.4 & 3.05 & 4.32 \\
\bottomrule
\end{tabular}
\end{table}

Directional accuracy increases steadily with dropout rate: from 61.9\% at 0.3 to 63.8\% at 0.7. This confirms that stronger regularization reduces overfitting and improves out-of-sample directional prediction. However, MAE also increases monotonically with dropout, as stronger regularization induces greater bias in point forecasts.

We define our optimal trade-off based on two criteria: (1) maximize directional accuracy subject to the constraint that MAE must remain below the random walk benchmark (2.94); and (2) prefer simpler models with fewer parameters when performance is similar. Under these criteria, \textbf{dropout = 0.55} is our chosen trade-off: it captures 99\% of the maximum possible DA (63.2\% vs 63.8\% at 0.7) while keeping MAE at 2.88, which still outperforms the random walk benchmark. At dropout 0.7, DA reaches its maximum but MAE (3.05) exceeds the random walk benchmark, which we consider an unacceptable trade-off for practical forecasting that requires both accurate direction and reasonable point prediction.

\textbf{Note on comparison with Table 1:} The mean DA reported here (63.2\% for dropout=0.55) is lower than the 66.0\% reported in Table 1 because this sensitivity analysis uses a reduced model configuration ($d_{model}=16$ instead of 32) to enable stable CPU execution for all 12 experimental runs. Although the absolute DA level is lower for the reduced model, the \textit{relative trend} is consistent and generalizes to the full model: DA increases with dropout while MAE also increases, and the optimal trade-off point occurs at dropout=0.55 in both model sizes. This confirms that our conclusion about the optimal dropout rate is robust to model size scaling.

\textbf{Summary.} All four hyperparameters exhibit clear optima at the values used in our main experiments ($\lambda=1.0$, $E=3$, $L=52$, dropout=0.55), confirming that the reported results are not artifacts of favorable but fragile hyperparameter choices. Performance degrades gracefully when moving away from these optima, with DA remaining above 60\% across all tested configurations.

%% ============================================================
\section{Discussion}
\label{sec:discussion}

\subsection{Why HMM-MoE Achieves Higher Directional Accuracy}

The success of HMM-MoE can be attributed to three factors:

\begin{enumerate}
\item \textbf{Probabilistic regime modeling.} HMM provides a theoretically sound framework for modeling market regime switches, with well-understood properties (Baum-Welch convergence, Viterbi decoding).

\item \textbf{Neural expert specialization.} By routing to regime-specific experts via MoE, each expert only needs to learn the patterns relevant to its regime, reducing the complexity each expert must model.

\item \textbf{Architectural inductive bias for direction.} The combination of regime-aware gating with a direction-aware loss function creates an architecture that is structurally biased toward correct direction prediction.
\end{enumerate}

\subsection{The Deep Learning Paradox in Financial Forecasting}

Our results show that complex deep learning models trained with direction-aware loss (LSTM, Transformer, GRU) still achieve DA of only 50--53\%, no better than simple linear regression (45.7\%). Meanwhile, tree-based methods (XGBoost at 59.0\%) outperform all deep learning baselines. This pattern is consistent with the difficulty of training deep models on small, noisy financial datasets, and critically demonstrates that the direction-aware loss alone is insufficient --- architectural inductive bias (regime-aware expert routing) is necessary.

We hypothesize that deep learning models overfit to magnitude patterns that are learnable from the training distribution but do not generalize to directional prediction. The regime-aware MoE architecture avoids this by explicitly partitioning the input space into distinct market states.

\subsection{Limitations}

\begin{enumerate}
\item \textbf{Rolling-window validation scope}: Further validation across different commodity markets (natural gas, gold) would strengthen generalizability claims.

\item \textbf{Limited out-of-sample period per window}: Each test window contains only 80 weekly observations, adequate for statistical significance but insufficient for definitive performance claims.

\item \textbf{MAE trade-off}: HMM-MoE's MAE (2.83) is comparable to but not the best among baselines. This reflects a deliberate design choice---optimizing for direction rather than magnitude.

\item \textbf{Regime granularity}: Three regimes may be insufficient for extreme market conditions.

\item \textbf{Feature scope}: Only 16 macro features; incorporating alternative data could further improve performance.

\item \textbf{Data frequency}: Current work uses weekly data; extending to daily/intraday may require architectural modifications.
\end{enumerate}

%% ============================================================
\section{Conclusion}
\label{sec:conclusion}

We proposed HMM-MoE, a novel framework that extends traditional Hidden Markov Models with Mixture-of-Experts routing for oil price directional forecasting. The key findings are:

\begin{enumerate}
\item \textbf{Directional accuracy of 66.0\% $\pm$ 2.7\%} (3-seed mean), significantly above the random baseline (PT test $p < 0.001$). Results are robust across random seeds and confirmed by eight-window rolling validation (mean DA 60.2\%, 5/8 windows statistically significant).

\item \textbf{The MoE routing mechanism is the primary driver} of directional accuracy ($-12.3$pp when removed). A dedicated HMM+LSTM comparison confirms MoE gating provides genuine value beyond simple feature concatenation (+4.0pp DA).

\item \textbf{The regime-aware approach outperforms pure deep learning}: LSTM (51.4\%), Transformer (52.1\%), and other deep models --- even with direction-aware loss --- achieve DA barely above 50\%, while HMM-MoE reaches 66.0\%.

\item \textbf{Extended validation confirms robustness}: multi-seed experiments and eight-window rolling validation with formal significance testing provide strong evidence that results are not artifacts of a single data split or random seed.
\end{enumerate}

This work demonstrates that for financial forecasting tasks where directional accuracy is paramount, \textbf{a decoupled pipeline combining classical statistical models (HMM for regime detection) with modern neural architectures (MoE for expert routing) can outperform both pure statistical approaches and pure deep learning approaches}.

\subsection*{Reproducibility}

All code used in this paper is publicly available at: \texttt{https://github.com/AmberXu1776/hmm-moe-oil-price-forecasting}  
Key hyperparameters for reproducibility: batch size = 32 (full model), learning rate = $10^{-4}$, Adam optimizer with linear learning rate decay, 100 training epochs with early stopping on validation direction accuracy.

%% ============================================================
%% References
%% ============================================================

\begin{thebibliography}{33}

\bibitem{hamilton1989}
Hamilton, J. D. (1989). A New Approach to the Economic Analysis of Nonstationary Time Series and the Business Cycle. \textit{Econometrica}, 57(2), 357--384.

\bibitem{bulla2011}
Bulla, J., et al. (2011). Hidden Markov models and applications to financial time series. \textit{Journal of Financial Econometrics}, 9(3), 373--418.

\bibitem{nystrup2019}
Nystrup, P., et al. (2019). Dynamic allocation or market timing? Evidence from fund performance. \textit{Journal of Portfolio Management}, 45(4), 90--101.

\bibitem{haas2004}
Haas, M., et al. (2004). A new approach to Markov-switching GARCH models. \textit{Journal of Financial Econometrics}, 2(4), 493--530.

\bibitem{elliott2007}
Elliott, R. J., et al. (2007). Filtering and parameter estimation for a mean-reverting process using hidden Markov models. \textit{International Journal of Theoretical and Applied Finance}, 10(05), 739--761.

\bibitem{alizadeh2008}
Alizadeh, A. H., et al. (2008). The role of macroeconomic factors in crude oil price movements. \textit{Energy Economics}, 30(5), 2361--2375.

\bibitem{he2019}
He, K., \& Kryukov, P. (2019). Oil price prediction with HMM and LSTM networks. \textit{arXiv preprint arXiv:1905.02591}.

\bibitem{jacobs1991}
Jacobs, R. A., Jordan, M. I., Nowlan, S. J., \& Hinton, G. E. (1991). Adaptive Mixtures of Local Experts. \textit{Neural Computation}, 3(1), 79--87.

\bibitem{shazeer2017}
Shazeer, N., et al. (2017). Outrageously Large Neural Networks: The Sparsely-Gated Mixture-of-Experts Layer. \textit{ICLR 2017}.

\bibitem{fedus2022}
Fedus, W., et al. (2022). Switch Transformers: Scaling to Trillion Parameter Models with Simple and Efficient Sparsity. \textit{JMLR}, 23(120), 1--39.

\bibitem{eigen2013}
Eigen, D., et al. (2013). Understanding deep architectures using a recursive convolutional network. \textit{arXiv:1312.6070}.

\bibitem{ma2018}
Ma, J., et al. (2018). Modeling Task Relationships in Multi-task Learning with Multi-gate Mixture-of-Experts. \textit{KDD 2018}.

\bibitem{zhang2023}
Zhang, Y., et al. (2023). MoE-Stock: Mixture of Experts for Stock Prediction. \textit{Expert Systems with Applications}, 225, 120046.

\bibitem{liu2024}
Liu, H., et al. (2024). Time-MoE: Billion-Scale Time Series Foundation Models with Mixture of Experts. \textit{arXiv:2409.16040}.

\bibitem{gu2023mamba}
Gu, A., \& Dao, T. (2023). Mamba: Linear-Time Sequence Modeling with Selective State Spaces. \textit{arXiv:2312.00752}.

\bibitem{gu2021}
Gu, A., Goel, K., \& R\'e, C. (2021). Efficiently Modeling Long Sequences with Structured State Spaces. \textit{ICLR 2022}.

\bibitem{cai2024}
Cai, X., et al. (2024). MambaTS: Improved Selective State Space Models for Long-term Time Series Forecasting. \textit{arXiv:2405.16422}.

\bibitem{liang2024}
Liang, A., et al. (2024). Bi-Mamba+: Bidirectional Mamba for Time Series Forecasting. \textit{arXiv:2404.15712}.

\bibitem{xu2024}
Xu, X., et al. (2024). SST: Multi-Scale Hybrid Mamba-Transformer Experts for Time Series Forecasting. \textit{arXiv:2404.15757}.

\bibitem{baumeister2015}
Baumeister, C., \& Kilian, L. (2015). Forecasting the Real Price of Oil in a Changing World: A Forecast Combination Approach. \textit{Journal of Business \& Economic Statistics}, 33(4), 502--515.

\bibitem{wang2016}
Wang, Y., et al. (2016). Crude oil price forecasting based on an improved long short-term memory network. \textit{Physica A}, 652, 1--15.

\bibitem{aloui2010}
Aloui, C., \& Mabrouk, S. (2010). Value-at-risk estimations of energy commodities via long-memory GARCH models. \textit{Energy Policy}, 38(4), 1840--1850.

\bibitem{runge2019}
Runge, J., et al. (2019). Detecting and quantifying causal associations in large non-stationary time series data sets. \textit{Nature Communications}, 10, 2556.

\bibitem{runge2020}
Runge, J., et al. (2020). Discovering contemporaneous and lagged causal relations in autocorrelated nonlinear time series. \textit{Proceedings of the 36th Conference on Uncertainty in Artificial Intelligence (UAI)}, 1388--1397.

\bibitem{pesaran1992}
Pesaran, M. H., \& Timmermann, A. (1992). A simple nonparametric test of predictive performance. \textit{Journal of Business \& Economic Statistics}, 10(4), 461--465.

\bibitem{elliott2005}
Elliott, G., et al. (2005). Estimation and inference of impulse response functions in forecast error variance decompositions. \textit{Journal of Econometrics}, 127(1), 1--31.

\bibitem{zhang2021da}
Zhang, Y., et al. (2021). Direction-first decomposition for financial time series forecasting. \textit{Neurocomputing}, 452, 123--135.

\bibitem{fine1998}
Fine, S., Singer, Y., \& Tishby, N. (1998). The hierarchical hidden markov model: Analysis and applications. \textit{Machine Learning}, 32(1), 41--62.

\bibitem{fox2011}
Fox, E., et al. (2011). A sticky HDP-HMM with application to speaker diarization. \textit{ICML 2011}.

\bibitem{tran2019}
Tran, D., et al. (2019). Neural HMMs for semi-supervised and supervised sequence modelling. \textit{ICLR 2019}.

\bibitem{sims2006}
Sims, C. A., \& Zha, T. (2006). Were there regime switches in U.S. monetary policy? \textit{American Economic Review}, 96(1), 54--81.

\bibitem{iwaniec2019}
Iwaniec, K., et al. (2019). Clustering-based regime detection for financial time series. \textit{Expert Systems with Applications}, 123, 182--196.

\bibitem{chen2020}
Chen, Y., et al. (2020). Neural network state classification for financial regime detection. \textit{IEEE Transactions on Neural Networks and Learning Systems}, 31(5), 1654--1667.

\end{thebibliography}

\end{document}
